\section{Introduction}
\label{sec:intro}
LLM-based agents are becoming a practical interface for GPU-kernel generation: they can draft an implementation, run it, inspect compiler or execution feedback, and iterate. Benchmarks such as KernelBench~\citep{kernelbench} and TritonBench~\citep{tritonbench} have made this setting measurable, while recent agents improve kernels through feedback, search, post-training, or skill memory~\citep{gpu_kernel_scientist,geak,sakana_cuda_engineer,kernel_smith,kernel_skill,adaexplore}. The setting is also increasingly multi-surface: agents may write CUDA, Triton, CuTe/CUTLASS-style templates, TileLang programs, or a restricted kernel DSL, each exposing different scheduling knobs and failure modes~\citep{tritonbench,geak,cutegen,tritonforge,hari2026improving,chu2025gpu,qu2026two}. However, prior studies document a gap between attempting an optimization and applying it effectively. KernelBench reports that providing optimization examples can encourage more aggressive transformations while increasing execution failures~\citep{kernelbench}. 

The problem is often not that the model lacks the names of useful optimizations. It can suggest tiling, shared-memory staging, vectorized loads, software pipelining, or hardware intrinsics. What is missing is the condition under which each optimization is valid and useful. Vectorized loads require alignment and valid tail handling; pipelining assumes staged data movement; matrix intrinsics require compatible tile shapes, layouts, and hardware support. A forward-search agent may therefore apply a familiar optimization to the wrong state, producing a compile error, a correctness failure, or a slower kernel. Pretraining captures much of the \emph{what} of optimization, but not reliably the \emph{when}.

We treat this missing ``when'' knowledge as the object to be learned. The memory item we need is not a free-form hint such as ``use vectorized loads'', but a skill in a narrow sense: a verified, code-anchored specification of an optimization intent, including where the intent acts on a code or IR surface, what carrier can guide implementation, the context in which it has been observed to work, the effect it tends to have, and the risks when its assumptions are violated. Existing kernel-agent memories expose parts of this signal through failure rules, profile-state knowledge bases, or expert-skill snippets~\citep{adaexplore,ksearch,kernelblaster,kernel_skill}, while post-training approaches may absorb it into model weights~\citep{kernel_smith,kevin_rl_kernel,cuda_l1}. 
These memories are useful, but when they are built from forward rollouts, their coverage is limited by the trajectories the agent was able to discover. The optimization experience they accumulate depends on what the backbone model can propose and successfully implement. When high performance requires unfamiliar combinations of hardware features and optimization steps, forward rollouts can remain concentrated on familiar implementation patterns, leaving the relevant code states and dependencies unexplored within a fixed budget. Expert kernels supply successful implementations directly, giving \sys a concrete basis for recovering optimization paths and applicability conditions that the agent may fail to discover on its own. Our goal is to externalize the missing ``when'' knowledge from expert artifacts instead.

We introduce \sys, a framework for inducing reusable skills from expert GPU-kernel lineages. The key observation is that expert kernels already contain the missing conditional knowledge, but only implicitly. 
Expert kernels are developed through a sequence of interdependent optimization decisions involving tiling, vectorization, parameter selection, and hardware features. However, the final code typically leaves the intermediate code states, optimization order, and conditions that make each step valid implicit.\sys recovers these decisions by walking an expert kernel backward through named, validation-gated deoptimization steps. Each accepted simplification is stored as its inverse forward transition, yielding a curriculum from a naive implementation toward the expert kernel; the forward transitions observed across experts become evidence for reusable skills. To our knowledge, \sys is the first LLM-kernel memory constructed by deoptimizing expert implementations rather than accumulating forward-rollout states.


% Original wording (kept for ease of restoration): To our knowledge, \sys is the first LLM-kernel memory construction method that automatically derives reusable optimization lineages by deoptimizing expert implementations, rather than accumulating only the states reached by naive forward search or organizing pre-existing human/agent-written skill snippets.

At reuse time, \sys retrieves skills by scope. 
Given a new target, it retrieves skills along the target case, surface language, platform, and already-applied optimizations. The target state need not be naive: \sys can start from any current kernel attempt and use the retrieved lineage to complete missing intents or repair mis-materialized ones. It then lets a downstream LLM materialize selected skills on the target code surface using their anchors and carriers.
This keeps the LLM in the role where it is useful---instantiating a validated code-anchored intent and repairing local mismatches under feedback---while the memory constrains it with validated preconditions and observed risks.

We evaluate \sys{} on two NVIDIA platforms, an H100 and an RTX PRO 6000, and on an Ascend 910B, inducing one library on each from five expert kernels. On those kernels \sys{} is correct on all ten workload--platform pairs at $1.12\times$ over the vendor references, where AdaExplore~\citep{adaexplore} and AccelOpt~\citep{accelopt} reach $0.25\times$ and $0.54\times$ and return nothing correct on four and two of them. An agent handed the same expert kernels as source recovers individual optimizations but not the combinations they belong to, and its kernels run up to $105\times$ slower. On five new operators per platform that contribute no skills, \sys{} reaches $1.73\times$ on H100 and $1.41\times$ on Ascend over that same agent; on the $100$ Level~2 problems of KernelBench~\citep{kernelbench}, ported to AscendC as the Fusion category of MultiKernelBench~\citep{multikernelbench}, mean speedup rises from $1.05\times$ to $1.32\times$ against \texttt{torch.compile} on NVIDIA and from $0.63\times$ to $2.09\times$ against eager PyTorch on Ascend. Conditioned skills therefore make expert optimization knowledge reusable well beyond the kernels it was extracted from.

\paragraph{Contributions.}
This paper introduces \sys{} and makes three technical contributions:
\begin{itemize}
    \item \textbf{Conditional skills that encode when to optimize.} We make the ``when'' of optimization explicit by pairing each optimization intent with the conditions under which it is valid and useful. Each skill includes code anchors and implementation carriers to guide materialization, together with effects, risks, verification evidence, and scope over workloads, languages, platforms, and prior optimization states.

    \item \textbf{Guided deoptimization for constructing conditional skills.} We walk expert kernels backward through validation-gated simplifications and reconstruct accepted steps as validated forward transitions, forming an executable curriculum from naive to expert. We then lift these transitions into conditional skills, using the intermediate code states, optimization dependencies, and validation outcomes as evidence for their applicability conditions.

    \item \textbf{Using conditional skills to guide when to optimize.} We retrieve conditional skills according to the target kernel's current state and context, using their applicability conditions to guide which optimizations to apply at each stage. Only skills admitted through held-out roundtrip validation are reused; a downstream LLM materializes the selected skills into optimized kernels under a compile/correctness/profile gate.
\end{itemize}
